Consider a money system consisting of $n$ coins. Each coin has a positive integer value. Your task is to calculate the number of distinct ordered ways you can produce a money sum $x$ using the available coins.
For example, if the coins are $\{2,3,5\}$ and the desired sum is $9$, there are $3$ ways:

\textbullet\ $2+2+5$
\textbullet\ $3+3+3$
\textbullet\ $2+2+2+3$

\vspace{0.3cm}\noindent\textbf{Input}

The first input line has two integers $n$ and $x$: the number of coins and the desired sum of money.
The second line has $n$ distinct integers $c_1,c_2,\dots,c_n$: the value of each coin.

\vspace{0.3cm}\noindent\textbf{Output}

Print one integer: the number of ways modulo $10^9+7$.

\vspace{0.3cm}\noindent\textbf{Constraints}

\textbullet\ $1 \le n \le 100$
\textbullet\ $1 \le x \le 10^6$
\textbullet\ $1 \le c_i \le 10^6$